\section{Estimation of the borrowing criterion}
\label{sec:supp-gain}

\begin{lemma}[Kernel moments and denominator]
\label{lem:denominator}
Under Condition~\ref{cond:kernel},
\begin{align}
\mu_h
&=h^df_X(x_0)\kappa_1\{1+o(1)\},
\label{eq:mu-expansion}\\
E K_h^2(X-x_0)
&=h^df_X(x_0)\kappa_2\{1+o(1)\}.
\label{eq:K2-expansion}
\end{align}
If \(I\) is a fold of size \(n_I\) with \(n_I/M\to\rho_I\in(0,1)\), then
\begin{equation}
\frac{\mathbb P_IK_h}{\mu_h}-1
=O_{\mathbb P}\{(n_Ih^d)^{-1/2}\}
=O_{\mathbb P}\{(Mh^d)^{-1/2}\}.
\label{eq:den-rate}
\end{equation}
\end{lemma}

\begin{proof}
Changing variables \(x=x_0+hu\),
\[
\mu_h=h^d\int K(u)f_X(x_0+hu)\,du.
\]
Continuity of \(f_X\), compact support of \(K\), and dominated convergence give \eqref{eq:mu-expansion}; the same argument gives \eqref{eq:K2-expansion}. Moreover,
\[
\operatorname{Var}(\mathbb P_IK_h)
=\frac1{n_I}\operatorname{Var}\{K_h(X-x_0)\}
=O(h^d/n_I).
\]
Since \(\mu_h\asymp h^d\), Chebyshev's inequality gives \eqref{eq:den-rate}.
\end{proof}

\begin{lemma}[Ratio linearization]
\label{lem:ratio}
Under Conditions~\ref{cond:sampling} and \ref{cond:kernel}, let \(a\) be fixed
conditional on the data outside an evaluation fold \(I\), with
\(n_I/M\to\rho_I\in(0,1)\). If the scaled numerator process is tight in
\(\ell^\infty(\mathcal Y)\), then
\begin{align}
&\sqrt{Mp_Mh^d}
\{\widehat F_{h,I}^{a}(y)-F_{h,0}(y)\}
\nonumber\\
&\qquad=
\sqrt{Mp_Mh^d}
(\mathbb P_I-P)
\left[
\frac{K_h}{\mu_h}
\{U_y(a)-F_{h,0}(y)\}
\right]
+o_{\mathbb P}(1)
\label{eq:ratio-linearization}
\end{align}
uniformly over \(y\in\mathcal Y\).
\end{lemma}

\begin{proof}
The numerator has mean zero, so
\[
\widehat F_{h,I}^{a}(y)-F_{h,0}(y)
=
\frac{\mathbb P_I[K_h\{U_y(a)-F_{h,0}(y)\}]}{\mathbb P_IK_h}.
\]
Write \(D_I=\mathbb P_IK_h/\mu_h-1\). By Lemma~\ref{lem:denominator},
\[
D_I=O_{\mathbb P}\{(n_Ih^d)^{-1/2}\}=o_{\mathbb P}(1).
\]
The assumed tightness and \(n_I/M\to\rho_I\) give
\[
\left\|
(\mathbb P_I-P)
\left[
\frac{K_h}{\mu_h}\{U_\cdot(a)-F_{h,0}(\cdot)\}
\right]
\right\|_{\mathcal Y}
=O_{\mathbb P}\{(Mp_Mh^d)^{-1/2}\}.
\]
Expanding \((1+D_I)^{-1}=1-D_I+o_{\mathbb P}(|D_I|)\), the scaled remainder is \(O_{\mathbb P}(|D_I|)=o_{\mathbb P}(1)\), proving \eqref{eq:ratio-linearization}.
\end{proof}

Conditional on a fitting fold, suppress its index and treat \(q\), \(m\), and \(D\) as fixed. Define
\[
\mathcal G_{M,h}=p_MG_{M,h},
\qquad
\mathcal H_{M,h}=p_MH_{M,h}.
\]

\begin{lemma}[Observable alignment identity]
\label{lem:observable-G}
Conditional on the candidate regressions,
\begin{equation}
G_{M,h}
=
E\left[
 v_{M,h}
 \left\langle
 D,\frac{R}{\pi_M}(Z-q)
 \right\rangle_\nu
\right].
\label{eq:supp-observable-G}
\end{equation}
\end{lemma}

\begin{proof}
Conditional on \(W\),
\[
E\left\{
\frac{R}{\pi_M(W)}(Z_y-q_y(X))
\,\Bigm|\,W
\right\}
=m_{0,y}(W)-q_y(X)=A_y(W).
\]
Integrating against \(D_y\,d\nu(y)\), multiplying by \(v_{M,h}\), and taking expectation proves \eqref{eq:supp-observable-G}.
\end{proof}

\begin{lemma}[Gain-estimation rates]
\label{lem:gain-rates}
Under Conditions~\ref{cond:sampling}--\ref{cond:candidates}, on a selection fold of size \(n_G\) with \(n_G/M\to\rho_G\in(0,1)\),
\begin{align}
\widehat{\mathcal G}_{h}-\mathcal G_{M,h}
&=O_{\mathbb P}\{(Mp_Mh^d)^{-1/2}\},
\label{eq:gain-G-rate}\\
\widehat{\mathcal H}_{h}-\mathcal H_{M,h}
&=O_{\mathbb P}\{(Mh^d)^{-1/2}\}.
\label{eq:gain-H-rate}
\end{align}
The same rates hold when \(\mu_h\) in \(v_{M,h}\) is replaced by the selection-fold kernel mean.
\end{lemma}

\begin{proof}
Let
\[
T_{G,M}
=
p_Mv_{M,h}
\left\langle
D,\frac{R}{\pi_M}(Z-q)
\right\rangle_\nu.
\]
Uniform boundedness of the paths and finiteness of \(\nu\) imply, on the support of \(K_h\),
\[
|T_{G,M}|
\le
\frac{CR}{p_Mh^d}.
\]
Since
\[
\Pr\{R=1,K_h(X-x_0)\ne0\}
=E[\pi_M(W)1\{K_h(X-x_0)\ne0\}]
=O(p_Mh^d),
\]
we obtain
\[
E(T_{G,M}^2)=O\{(p_Mh^d)^{-1}\}.
\]
An empirical mean over \(n_G\asymp M\) observations therefore has standard deviation \(O\{(Mp_Mh^d)^{-1/2}\}\).

For
\[
T_{H,M}=p_Mv_{M,h}\|D\|_\nu^2,
\]
we have \(|T_{H,M}|\le Ch^{-d}\) on the kernel support, so
\[
E(T_{H,M}^2)=O(h^{-d}).
\]
This gives \eqref{eq:gain-H-rate}. Lemma~\ref{lem:denominator} shows that replacing \(\mu_h^{-2}\) by its selection-fold counterpart multiplies each criterion term by
\[
1+O_{\mathbb P}\{(Mh^d)^{-1/2}\}.
\]
The normalized population criteria are \(O(1)\), and \((Mh^d)^{-1/2}\le(Mp_Mh^d)^{-1/2}\), so this substitution is absorbed by the displayed rates.
\end{proof}

\begin{lemma}[Quadratic regret with numerical stabilization]
\label{lem:regret}
Under the conditions of Lemma~\ref{lem:gain-rates} and the stabilizer clauses
of Condition~\ref{cond:candidates}, let
\[
\mathcal Q_{M,h}(\lambda)
=-2\lambda\mathcal G_{M,h}+\lambda^2\mathcal H_{M,h}
\]
and define the stabilized empirical criterion
\[
\widehat{\mathcal Q}_{h}(\lambda)
=
-2\lambda\widehat{\mathcal G}_{h}
+
\lambda^2\max\{\widehat{\mathcal H}_{h},\varepsilon_M\}.
\]
Let \(\widehat\lambda_h\) and \(\lambda_h^{\mathrm{or}}\) minimize these criteria over \([0,1]\), and put
\[
r_{M,h}=(Mp_Mh^d)^{-1/2}.
\]
Then
\begin{align}
\sup_{\lambda\in[0,1]}
|\widehat{\mathcal Q}_{h}(\lambda)-\mathcal Q_{M,h}(\lambda)|
&=O_{\mathbb P}(r_{M,h}+\varepsilon_M),
\label{eq:Q-uniform}\\
\mathcal Q_{M,h}(\widehat\lambda_h)
-
\mathcal Q_{M,h}(\lambda_h^{\mathrm{or}})
&=O_{\mathbb P}(r_{M,h}+\varepsilon_M).
\label{eq:Q-regret}
\end{align}
If \(\mathcal H_{M,h}\to H_\infty>0\), then
\begin{equation}
\widehat\lambda_h-\lambda_h^{\mathrm{or}}
=O_{\mathbb P}(r_{M,h}).
\label{eq:lambda-direct-rate}
\end{equation}
If, in addition, \(\lambda_h^{\mathrm{or}}\in[\delta,1-\delta]\) for some \(\delta>0\) and all sufficiently large \(M\), then
\begin{equation}
\mathcal Q_{M,h}(\widehat\lambda_h)
-
\mathcal Q_{M,h}(\lambda_h^{\mathrm{or}})
=O_{\mathbb P}(r_{M,h}^2).
\label{eq:Q-interior-regret}
\end{equation}
If \(\mathcal H_{M,h}\to0\) and \(\sup_M\mathcal J_{M,h}(0)<\infty\), then
\begin{equation}
\sup_{\lambda\in[0,1]}
|\mathcal Q_{M,h}(\lambda)-\mathcal Q_{M,h}(0)|
=o(1).
\label{eq:flat-criterion}
\end{equation}
\end{lemma}

\begin{proof}
For \(0\le\lambda\le1\),
\[
\begin{aligned}
&|\widehat{\mathcal Q}_{h}(\lambda)-\mathcal Q_{M,h}(\lambda)|\\
&\qquad\le
2|\widehat{\mathcal G}_{h}-\mathcal G_{M,h}|
+
|\widehat{\mathcal H}_{h}-\mathcal H_{M,h}|
+
\varepsilon_M.
\end{aligned}
\]
Lemma~\ref{lem:gain-rates} proves \eqref{eq:Q-uniform}. Empirical optimality gives
\[
\mathcal Q_{M,h}(\widehat\lambda_h)
-
\mathcal Q_{M,h}(\lambda_h^{\mathrm{or}})
\le
2\sup_{\lambda\in[0,1]}
|\widehat{\mathcal Q}_{h}(\lambda)-\mathcal Q_{M,h}(\lambda)|,
\]
which proves \eqref{eq:Q-regret}.

Suppose \(\mathcal H_{M,h}\to H_\infty>0\). Then \(\widehat{\mathcal H}_{h}\to_{\mathbb P}H_\infty\), so the stabilizer is inactive with probability tending to one. On that event,
\[
\widehat\lambda_h
=
\Pi_{[0,1]}
\left(
\frac{\widehat{\mathcal G}_{h}}{\widehat{\mathcal H}_{h}}
\right).
\]
Projection is nonexpansive and the ratio map is Lipschitz in a neighborhood of \(H_\infty\), which gives \eqref{eq:lambda-direct-rate}.

If \(\lambda_h^{\mathrm{or}}\in[\delta,1-\delta]\), then the population minimizer is interior and \(\mathcal G_{M,h}=\mathcal H_{M,h}\lambda_h^{\mathrm{or}}\). Equation \eqref{eq:lambda-direct-rate} places \(\widehat\lambda_h\) in \((0,1)\) with probability tending to one. The exact quadratic identity yields
\[
\mathcal Q_{M,h}(\widehat\lambda_h)
-
\mathcal Q_{M,h}(\lambda_h^{\mathrm{or}})
=
\mathcal H_{M,h}
(\widehat\lambda_h-\lambda_h^{\mathrm{or}})^2,
\]
which is \(O_{\mathbb P}(r_{M,h}^2)\).

Now suppose \(\mathcal H_{M,h}\to0\). Weighted Cauchy--Schwarz gives
\[
|\mathcal G_{M,h}|^2
\le
\mathcal J_{M,h}(0)\mathcal H_{M,h},
\]
so \(\mathcal G_{M,h}\to0\). Uniformly over \(\lambda\in[0,1]\),
\[
|\mathcal Q_{M,h}(\lambda)-\mathcal Q_{M,h}(0)|
\le2|\mathcal G_{M,h}|+\mathcal H_{M,h}
=o(1),
\]
which proves \eqref{eq:flat-criterion}.
\end{proof}

\begin{lemma}[Primitive norm comparison]
\label{lem:norm-comparison}
Under Conditions~\ref{cond:sampling} and \ref{cond:kernel}, suppose \(\nu\)
has a density bounded below on \(\mathcal Y\), and every path \(D_M(W)\) is
uniformly \(\alpha\)-H\"older with constant \(L_D\):
\[
|D_{M,z}(W)-D_{M,y}(W)|
\le L_D|z-y|^\alpha.
\]
Then Condition~\ref{cond:path-compatibility} holds with
\[
\omega(t)=Ct^{2\alpha/(2\alpha+1)}.
\]
\end{lemma}

\begin{proof}
Let \(f\) be an \(\alpha\)-H\"older function with constant \(L_D\), and set \(a=\|f\|_\infty\). On an interval of length proportional to \(a^{1/\alpha}\) around a maximizer of \(|f|\), one has \(|f|\ge a/2\). Since the density of \(\nu\) is bounded below,
\[
\|f\|_\nu^2\ge Ca^{2+1/\alpha},
\qquad
\|f\|_\infty^2
\le C\{\|f\|_\nu^2\}^{\gamma},
\qquad
\gamma=\frac{2\alpha}{2\alpha+1}.
\]
Let \(w_{M,h}(W)=p_Mv_{M,h}(W)\). Then
\[
\mathcal D_{M,h}^2
\le E\{w_{M,h}\|D_M\|_\infty^2\}
\le C E\{w_{M,h}(\|D_M\|_\nu^2)^\gamma\}.
\]
Because \(\gamma\in(0,1)\), Jensen's inequality under the probability measure with density \(w_{M,h}/Ew_{M,h}\) gives
\[
E\{w_{M,h}(\|D_M\|_\nu^2)^\gamma\}
\le
(Ew_{M,h})^{1-\gamma}
\{E(w_{M,h}\|D_M\|_\nu^2)\}^{\gamma}.
\]
Conditions~\ref{cond:sampling}--\ref{cond:kernel} imply \(Ew_{M,h}=O(1)\), while the last expectation is \(\mathcal H_{M,h}\). Hence
\[
\mathcal D_{M,h}^2
\le C\mathcal H_{M,h}^{2\alpha/(2\alpha+1)}+o(1),
\]
which proves the claim.
\end{proof}
